\chapter{Las sucesiones monótonas acotadas son de Cauchy}

\section{Subsucesiones iteradas y órbitas}

Comencemos esta lectura con un ejercicio fundamental que iluminará la
profunda conexión entre la iteración y el crecimiento monótono.

Supongamos que tenemos una sucesión de números naturales,
$\sigma : \mathbb{N} \to \mathbb{N}$, y lo único que sabemos de ella es
que siempre supera a la función identidad:
\[
  h_\sigma : \forall n,\; n < \sigma(n)
\]

Esto no significa que $\sigma(n)$ sea en sí misma estrictamente creciente
(lo que llamamos \emph{Subsucc}). La sucesión puede saltar de un lado a
otro, siempre que su gráfica permanezca por encima de la recta $y = x$.

Pero debería ser ``intuitivamente claro'' a partir de $h_\sigma$ que
$\sigma$ acaba creciendo indefinidamente, haciéndose cada vez más grande,
aunque no de manera monótona. Es decir, debería existir alguna forma de
``acelerar'' $\sigma$ para convertirla en una Subsucc. El único problema
es: ¿cómo se hace esto en la práctica?

\subsection{La idea clave: órbitas}

La idea clave es la de una \emph{órbita}. En astronomía, podemos imaginar
que observamos el cielo noche tras noche e intentamos rastrear la posición
de, por ejemplo, Júpiter frente a las estrellas ``fijas'' (la esfera
celeste). Iniciamos nuestras observaciones con Júpiter teniendo cierto
estado en el ``espacio de fases'' (posición, velocidad) $x_0$; sea $T$ la
función que aplica la dinámica newtoniana durante un día, de modo que
$T(x_0)$ es el nuevo estado de Júpiter mañana, moviéndose según las leyes
de Newton y la gravedad. Entonces $T(T(x_0))$ es el estado después de dos
días, y así sucesivamente. La órbita completa de Júpiter en el tiempo es
la sucesión:
\[
  x_0,\; T(x_0),\; T(T(x_0)),\; T(T(T(x_0))),\; \ldots
\]

En matemáticas, si tenemos cualquier función $f : X \to X$ que va de un
espacio abstracto $X$ a sí mismo, y empezamos con algún punto base
$x_0 : X$, escribiremos $f^{[n]}(x_0)$ para denotar $f$ iterada $n$ veces
aplicada a $x_0$. La sucesión $n \mapsto f^{[n]}(x_0)$ se denomina
\emph{órbita de $x_0$ bajo la acción de $f$}.

\subsection{Aplicación a nuestro problema}

¿Cómo nos ayuda esto en nuestra situación actual? Podríamos empezar con
cualquier punto base $n_0 : \mathbb{N}$, y sabemos por $h_\sigma$
especializado a $n = n_0$ que $n_0 < \sigma(n_0)$, pero no tenemos idea
de cuán grande es $\sigma(n_0)$; podría ser enorme. Entonces, ¿cómo
garantizamos que el siguiente término supere a $\sigma(n_0)$? (¿Quiere
pensarlo un momento antes de seguir leyendo?)

Dado lo que discutimos antes, esperamos que se vea de inmediato que: si
especializamos $h_\sigma$ a $n = \sigma(n_0)$, obtenemos:
$\sigma(n_0) < \sigma(\sigma(n_0))$. Así queda claro: la manera de
obtener términos cada vez más grandes de la sucesión $\sigma$ es
\emph{tomar la órbita}.

\subsection{Nuevas herramientas}

\subsubsection{Iteración de funciones: \texttt{succ\_iterate}}

Aunque $\sigma^{[k]}(\sigma(n)) = \sigma^{[k+1]}(n)$ es verdad por
definición, se requiere un argumento por inducción para mostrar que si en
lugar de añadir una $\sigma$ a la derecha la añadimos a la izquierda:
\[
  \sigma(\sigma^{[k]}(n)) = \sigma^{[k+1]}(n)
\]
Omitiremos la demostración de ese argumento y les proporcionamos el
teorema \texttt{succ\_iterate}.

\subsubsection{Subsucesión a partir del sucesor: \texttt{subseq\_of\_succ}}

Para demostrar que $\sigma$ es una Subsucc, cuya definición habla de todos
los $i < j$, basta hacerlo paso a paso. El teorema \texttt{subseq\_of\_succ}
dice que es suficiente mostrar que $\sigma(n) < \sigma(n+1)$ se cumple
para todo $n$ para concluir que $\sigma$ es una Subsucc. Se puede aplicar
este hecho para reducir la demostración a mostrar que $\sigma$ crece de
$n$ a $n+1$.

\subsubsection{Táctica: \texttt{show}}

La táctica \texttt{show} tiene la sintaxis \texttt{show hecho by prueba}.
Por ejemplo, si se desea reescribir usando el hecho de que
$\sigma(\sigma^{[n]}(n_0)) = \sigma^{[n+1]}(n_0)$ sin una declaración
\texttt{have} separada, se puede escribir:

\begin{minted}{lean4}
rewrite [ show σ (σ^[ n ] n₀) = σ^[ n+1] n₀ by apply succ_iterate ]
\end{minted}

\subsection{El enunciado matemático}

\begin{theorem}[Subsucc del Iterado]
Si una sucesión $\sigma : \mathbb{N} \to \mathbb{N}$ crece más rápido que
la identidad, $n < \sigma(n)$, entonces la órbita de cualquier punto base
$n_0 : \mathbb{N}$ bajo $\sigma$ --- es decir, la sucesión
$n_0, \sigma(n_0), \sigma^{[2]}(n_0), \ldots$ --- es una Subsucc, esto
es, es estrictamente creciente.
\end{theorem}

\subsection{Estrategia de la demostración}

La demostración usa el teorema \texttt{subseq\_of\_succ} para reducir el
problema a mostrar que los términos consecutivos en la órbita están
estrictamente ordenados. Específicamente, necesitamos mostrar:
\[
  \sigma^{[n]}(n_0) < \sigma^{[n+1]}(n_0)
\]
La idea clave es que $\sigma^{[n+1]}(n_0) = \sigma(\sigma^{[n]}(n_0))$
por el teorema \texttt{succ\_iterate}, y podemos aplicar la hipótesis
$h_\sigma$ al punto $\sigma^{[n]}(n_0)$ para obtener la desigualdad
deseada.

\subsection{Demostración en Lean}

\begin{minted}{lean4}
Statement Subseq_of_Iterate (σ : ℕ → ℕ) (hσ : ∀ n, n < σ n) (n₀ : ℕ) :
    Subseq (fun n ↦ σ^[n] n₀) := by
  apply subseq_of_succ
  intro n
  specialize hσ (σ^[n] n₀)
  have f : σ (σ^[n] n₀) = σ^[n+1] n₀ := by apply succ_iterate
  rewrite [f] at hσ
  apply hσ
\end{minted}

\subsection{Demostración en lenguaje natural}

\begin{proof}
Usamos el teorema \texttt{subseq\_of\_succ}, que reduce demostrar que una
sucesión es estrictamente creciente a mostrar que los términos consecutivos
satisfacen $f(n) < f(n+1)$.

Para la sucesión de órbita $n \mapsto \sigma^{[n]}(n_0)$, necesitamos
mostrar:
\[
  \sigma^{[n]}(n_0) < \sigma^{[n+1]}(n_0)
\]

Por el teorema \texttt{succ\_iterate}:
\[
  \sigma^{[n+1]}(n_0) = \sigma(\sigma^{[n]}(n_0))
\]

Ahora aplicamos la hipótesis $h_\sigma$ al punto $\sigma^{[n]}(n_0)$:
\begin{align*}
  \sigma^{[n]}(n_0) &< \sigma(\sigma^{[n]}(n_0)) \\
               &= \sigma^{[n+1]}(n_0)
\end{align*}

Esto completa la demostración.
\end{proof}

\subsection{Por qué esto importa}

Este resultado establece que incluso las sucesiones no monótonas pueden
tener subsucesiones monótonas cuando satisfacen ciertas condiciones de
crecimiento. La construcción de la órbita proporciona una manera canónica
de extraer comportamiento monótono de sucesiones que crecen más rápido que
la identidad. Esta técnica será crucial para el teorema principal sobre
sucesiones monótonas acotadas.

\section{Elección mejorada: de la existencia a la estructura}

¿Conoce la ``conjetura de los primos gemelos''? Un número $n : \mathbb{N}$
se llama \emph{Primo Gemelo} si tanto $n$ como $n+2$ son números primos.
Llamemos a esta propiedad $p : \mathbb{N} \to \mathsf{Prop}$.

Así, $p$ es una función que toma un número natural $n$ y devuelve
$p(n) = \text{Sí}/\text{No}$, dependiendo de si $n$ es un primo gemelo.

La conjetura de los primos gemelos afirma que hay infinitos primos gemelos;
es decir, para cualquier cota $N$, por grande que sea, siempre existe al
menos un $n > N$ para el cual $p(n)$ se cumple. Enunciaríamos la conjetura
así:
\[
  h : \forall N,\; \exists\, n > N,\; p(n)
\]

Supongamos ahora que tenemos una hipótesis como esta, para alguna propiedad
abstracta $p$ (si se prefiere, puede seguir pensando en $p(n)$ como la
prueba de si $n$ es un primo gemelo). Dado que existen $n$'s
arbitrariamente grandes para los cuales $p(n)$ se cumple, ¿cómo obtenemos
una subsucesión $\sigma : \mathbb{N} \to \mathbb{N}$, tal que a lo largo de
la subsucesión, $p(\sigma(n))$ se cumpla para todo $n$?

\subsection{Interpretar la existencia como funciones}

La idea es que debería interpretarse $h$ como una colección de funciones.
Dado un número natural $N$, la hipótesis $h$ nos producirá un $n$, pero
ese $n$ es función de $N$, así que deberíamos escribir $n = n(N)$. ¡Pero
eso no es todo! La hipótesis $h$ también contiene una demostración de que
$n(N) > N$ para todo $N$. ¿Se ve por qué el enunciado de $h$ implica la
existencia de tal función? Y por último, $h$ también nos da una
demostración del hecho de que, para todo $N$, $p(n(N))$ se cumple. ¡Así
que $n(N)$ es (casi) nuestra sucesión buscada! (Como resulta extraño
escribir $n(N)$, renombramos $n$ como $\tau$, para poder escribir
$\tau(N)$.)

La forma de pasar de $h$ a estas sucesiones en Lean es invocar una táctica
familiar: \texttt{choose}. Si se escribe:

\begin{minted}{lean4}
choose τ hτBnd hτP using h
\end{minted}

entonces Lean añadirá al estado del objetivo:

\begin{minted}{lean4}
τ : ℕ → ℕ
hτBnd : ∀ N, τ N > N
hτP  : ∀ N, p (τ N)
\end{minted}

¿No es sorprendente? Ahora debería poderse usar la idea de la sección
anterior para construir la subsucesión $\sigma$ deseada, aunque habrá que
pensar cuidadosamente en cómo hacerlo exactamente\ldots

\subsection{De la existencia a las subsucesiones}

Dado este contexto, nuestro objetivo es construir una subsucesión
$\sigma : \mathbb{N} \to \mathbb{N}$ que sea estrictamente creciente y
preserve la propiedad $p$. Es decir, queremos:
\begin{itemize}[nosep]
  \item $\sigma$ es una Subsucc (estrictamente creciente)
  \item $p(\sigma(n))$ se cumple para todo $n$
\end{itemize}

El desafío es que la función $\tau$ extraída por \texttt{choose} satisface
la propiedad pero puede no ser monótona. Sin embargo, podemos aplicar la
construcción de la órbita de la sección anterior para imponer estructura
monótona.

\subsection{El enunciado matemático}

\begin{theorem}[Elección mejorada]
Para cualquier propiedad $p : \mathbb{N} \to \mathsf{Prop}$, si existen
números naturales arbitrariamente grandes que satisfacen $p$, entonces
existe una subsucesión estrictamente creciente a lo largo de la cual $p$
siempre se cumple.

Formalmente: si $\forall N,\, \exists\, n > N,\, p(n)$, entonces
$\exists\, \sigma,\; \mathit{Subsucc}(\sigma) \land \forall n,\, p(\sigma(n))$.
\end{theorem}

\subsection{Estrategia de la demostración}

El desafío es que la función $\tau$ extraída por \texttt{choose} satisface
la propiedad pero puede no ser monótona. Sin embargo, podemos aplicar la
construcción de la órbita de la sección anterior para imponer estructura
monótona.

La demostración sigue un proceso natural en dos pasos:

\textbf{Paso 1: Extraer la función oculta.} Usar \texttt{choose} para
convertir el enunciado de existencia en funciones concretas como se mostró
arriba.

\textbf{Paso 2: Aplicar la construcción de la órbita.} Definir
$\sigma(n) = \tau^{[n]}(\tau(0))$. La construcción de la órbita garantiza
que $\sigma$ es estrictamente creciente (por el resultado de la Sección
1), mientras que la elección del punto inicial $\tau(0)$ y la iteración
preservan la propiedad $p$.

\subsection{Demostración en Lean}

\begin{minted}{lean4}
Statement (p : ℕ → Prop) (h : ∀ N, ∃ n > N, p n) :
    ∃ σ, Subseq σ ∧ ∀ n, p (σ n) := by
  choose τ hτBnd hτP using h
  let σ : ℕ → ℕ := fun n ↦ τ^[n] (τ 0)
  use σ
  split_ands
  · apply Subseq_of_Iterate τ hτBnd (τ 0)
  · intro n
    change p (τ^[n] (τ 0))
    rewrite [← show τ (τ^[n-1] (τ 0)) = τ^[n] (τ 0) by apply succ_iterate]
    apply hτP (τ^[n-1] (τ 0))
\end{minted}

\subsection{Demostración en lenguaje natural}

\begin{proof}
Extraemos de la hipótesis $h$ una función $\tau : \mathbb{N} \to \mathbb{N}$
junto con demostraciones de que $\tau(N) > N$ y $p(\tau(N))$ para todo $N$.

A continuación, definimos nuestra subsucesión deseada como
$\sigma(n) = \tau^{[n]}(\tau(0))$, la órbita de $\tau$ comenzando desde
el punto $\tau(0)$.

Para verificar que $\sigma$ es una subsucesión, aplicamos el resultado de
la Sección 1: como $\tau(N) > N$ para todo $N$, la órbita
$n \mapsto \tau^{[n]}(\tau(0))$ es estrictamente creciente.

Para verificar que $p(\sigma(n))$ se cumple para todo $n$, usamos el
hecho de que $\sigma(n) = \tau^{[n]}(\tau(0))$. Por el teorema
\texttt{succ\_iterate}, esto es igual a $\tau(\tau^{[n-1]}(\tau(0)))$
cuando $n > 0$. Como $p(\tau(M))$ se cumple para cualquier $M$, y
$\tau^{[n-1]}(\tau(0))$ es tal $M$, tenemos $p(\sigma(n))$ según se
requería. Para $n = 0$, tenemos $\sigma(0) = \tau(0)$, y $p(\tau(0))$ se
cumple por nuestra extracción.
\end{proof}

\subsection{Por qué este puente importa}

Este resultado demuestra el patrón fundamental de transformar enunciados
de existencia abstractos en objetos concretos y estructurados. La técnica
de usar \texttt{choose} para extraer funciones ocultas y luego aplicar
métodos de construcción para imponer estructura adicional aparece a lo
largo de las matemáticas formales.

Más específicamente, este resultado proporciona exactamente la herramienta
que necesitamos para el teorema principal. Cuando suponemos que una
sucesión monótona acotada no es de Cauchy, obtenemos un enunciado de
existencia sobre brechas persistentes. Este resultado intermedio muestra
cómo extraer de tales enunciados las subsucesiones concretas necesarias
para nuestro argumento de contradicción.

La combinación fluida de intuición de teoría de números (conjetura de los
primos gemelos) con técnicas analíticas (construcción de órbitas) ilustra
cómo diferentes áreas de las matemáticas se refuerzan y enriquecen
mutuamente.

\section{Las sucesiones monótonas y acotadas son de Cauchy}

El teorema fundamental que conecta monotonicidad, acotamiento y
convergencia establece que toda sucesión monótona acotada es de Cauchy (y
por lo tanto convergente). Este resultado es central en el análisis real y
proporciona una poderosa herramienta para establecer la convergencia sin
encontrar explícitamente el límite.

La intuición puede parecer clara: una sucesión monótona que está acotada
no puede tener brechas persistentes, porque tales brechas acabarían
haciendo que la sucesión superara su cota. Pero hacer que los detalles
sean correctos requerirá algo de trabajo. Posponemos a la siguiente
sección la demostración de un ``lema auxiliar'' y pasamos directamente al
argumento principal.

\subsection{Abstracción y generalidad}

Sea $a : \mathbb{N} \to X$ una sucesión. Un momento, ¿qué es $X$ aquí?
Bueno, estamos trabajando hacia la construcción de los números reales, así
que es mejor no presuponer su existencia; quizás $X$ debería ser los
racionales. Por otro lado, una vez que demostremos el hecho de que las
sucesiones monótonas acotadas son de Cauchy, quizás querramos usar ese
hecho para sucesiones de valores reales. ¿Así que necesitaríamos demostrar
el teorema dos veces? Por supuesto que no; ¡esa es la belleza de la
abstracción!

Estableceremos $X$ como un ``Tipo'' abstracto, pero supondremos cosas
sobre él como que existe un orden lineal (para poder decir $x \leq y$),
y una norma (para poder decir $|x|$), y que podemos sumar, restar,
multiplicar y dividir elementos de $X$ obteniendo elementos de $X$ (que
$X$ es un ``campo''). Haremos solo las suposiciones suficientes para que
la demostración funcione. Y luego, si queremos aplicar este teorema
general a una sucesión racional $a : \mathbb{N} \to \mathbb{Q}$ o una
sucesión real $a : \mathbb{N} \to \mathbb{R}$, Lean inferirá
automáticamente que todas las condiciones necesarias sobre $X$ se
satisfacen en esos casos.

\subsection{Nuevas herramientas}

\subsubsection{Definición: monótona}

Una sucesión $a : \mathbb{N} \to X$ (donde $X$ es un tipo ordenado) se
dice que es \emph{monótona} si $a(n) \leq a(m)$ siempre que $n \leq m$.

\subsubsection{Teorema: monotonía del sucesor}

Para demostrar la monotonicidad, basta hacerlo paso a paso; es decir, si
$a(m) \leq a(m+1)$ se cumple para todo $m$, entonces $a$ es monótona.

\subsubsection{Táctica: \texttt{push\_neg}}

La negación de $\forall$ es $\exists$, y viceversa. Para propagar una
cadena de negaciones, se escribe \texttt{push\_neg}.

\subsection{El enunciado matemático}

\begin{theorem}[\texttt{IsCauchyOfMonotoneBdd}]
Sea $a : \mathbb{N} \to X$ una sucesión monótona acotada superiormente por
algún $M : X$. Entonces $a$ es de Cauchy.
\end{theorem}

\subsection{Estrategia de la demostración}

¿Qué podría salir mal? Supongamos (por contradicción) que la sucesión $a$
no es de Cauchy. Esto significará que existen ``brechas'' de tamaño
$\varepsilon$ arbitrariamente tardías en la sucesión. El teorema de la
brecha iterada (véase la siguiente sección) muestra entonces que estas
brechas se acumulan sin límite, contradiciendo el supuesto de acotamiento.

Los pasos iniciales ``obvios'' son los siguientes:

\begin{enumerate}
  \item Suponer por contradicción la hipótesis $h$ de que $a$ no es de
    Cauchy. Entonces existe algún $\varepsilon > 0$ tal que para cualquier
    $N$, podemos encontrar $m \geq n \geq N$ con
    $|a(m) - a(n)| \geq \varepsilon$. (La táctica \texttt{push\_neg}
    resultará útil aquí, y quizás se quiera usar \texttt{change at h}
    para que Lean vea a través de la definición de \texttt{IsCauchy}.)

  \item Una vez que se elige este $\varepsilon$, se tiene el enunciado:
    $h : \forall N,\, \exists\, n \geq N,\, m \geq N,\, |a(m) - a(n)| \geq \varepsilon$.
    Piénsese detenidamente: $n$ es en realidad una función de $N$, es
    decir, $n = n(N)$; lo mismo con $m = m(N)$. Y dado un $N$, también se
    obtiene una prueba de que $n(N) \geq N$, una prueba de que
    $m(N) \geq n(N)$, y una prueba de que
    $|a(m(N)) - a(n(N))| \geq \varepsilon$.

  \item Ahora podemos usar \texttt{choose} con $h$ para extraer estas
    aplicaciones $n(N)$ y $m(N)$, excepto que la situación queda más
    clara si las renombramos como $\tau(N)$ y $\sigma(N)$. El hecho de
    que $\tau(N) \geq N$ puede llamarse \texttt{hτ} y el hecho de que
    $\sigma(N) \geq \tau(N)$ es \texttt{hσ}. El hecho de que
    $|a(\sigma(N)) - a(\tau(N))| \geq \varepsilon$ podemos llamarlo
    \texttt{hgap}.
\end{enumerate}

¿Y ahora qué? La intuición debería ser al menos algo clara: la sucesión
$a$ crece al menos en $\varepsilon$ de vez en cuando, y es creciente, así
que después de suficientes esos ``saltos'' hacia arriba de $\varepsilon$,
superaremos la cota superior declarada $M$. Pero, ¿cómo lo implementamos
en la práctica?

Bueno, empecemos con $N = 0$. Entonces obtenemos algún valor
$0 \leq \tau(0) \leq \sigma(0)$ para el cual
$|a(\sigma(0)) - a(\tau(0))| \geq \varepsilon$. Por supuesto podemos
eliminar los valores absolutos, dado que $a$ es monótona. ¿Cómo
continuamos?

\subsection{Lema auxiliar pospuesto}

Una vez que obtenemos cierta cantidad $\varepsilon$ de crecimiento,
necesitamos iterarlo para obtener $k \cdot \varepsilon$ de crecimiento,
para cualquier $k$:

\begin{theorem}[\texttt{IterateGap}]
Dados $(a : \mathbb{N} \to X)$, $(ha : \mathit{Monotone}\, a)$,
$(\varepsilon : X)$, $(\varepsilon_{\mathit{pos}} : \varepsilon > 0)$,
$(\tau : \mathbb{N} \to \mathbb{N})$,
$(h_\tau : \forall n,\, \tau(n) \geq n)$,
$(\sigma : \mathbb{N} \to \mathbb{N})$,
$(h_\sigma : \forall n,\, \sigma(n) \geq \tau(n))$,
$(h_{\mathit{gap}} : \forall n,\, \varepsilon \leq |a(\sigma(n)) - a(\tau(n))|)$:
\[
  \forall k,\quad k \cdot \varepsilon \leq a(\sigma^{[k]}(0)) - a(0)
\]

Es decir, la órbita $\sigma^{[k]}(0)$ es exactamente la subsucesión a lo
largo de la cual se garantiza que $a$ crezca al menos $k \times
\varepsilon$.
\end{theorem}

\subsection{Demostración en Lean}

\begin{minted}{lean4}
Statement IsCauchyOfMonotoneBdd {X : Type*}
    [NormedField X] [LinearOrder X] [IsStrictOrderedRing X]
    [FloorSemiring X] (a : ℕ → X) (M : X)
    (hM : ∀ n, a n ≤ M) (ha : Monotone a)
    : IsCauchy a := by
  intro ε hε
  by_contra h
  push_neg at h
  choose τ hτ σ hσ hgap using h
  have f1 : ∀ k, k * ε ≤ a (σ^[k] 0) - a 0 :=
    IterateGap a ha ε hε τ hτ σ hσ hgap
  let k : ℕ := ⌈(M - a 0) / ε⌉₊ + 1
  have hk' : (M - a 0) / ε ≤ ⌈(M - a 0) / ε⌉₊ := by bound
  have hk : (M - a 0) / ε < k := by
    change (M - a 0) / ε < (⌈(M - a 0) / ε⌉₊ + 1 : ℕ)
    push_cast; linarith [hk']
  specialize f1 k
  specialize hM (σ^[k] 0)
  have f2 : (M - a 0) < k * ε := by
    field_simp at hk
    rewrite [show k * ε = ε * k by ring_nf]
    apply hk
  linarith [f1, f2, hM]
\end{minted}

\subsection{Demostración en lenguaje natural}

\begin{proof}
Procedemos por contradicción. Supongamos que $a$ no es de Cauchy. Entonces
existe algún $\varepsilon > 0$ tal que para todo $N \in \mathbb{N}$,
podemos encontrar índices $m \geq n \geq N$ con
$|a(m) - a(n)| \geq \varepsilon$.

Eso significa que podemos construir funciones $\tau, \sigma : \mathbb{N}
\to \mathbb{N}$ tales que:
\begin{itemize}[nosep]
  \item $\tau(k) \geq k$ para todo $k$
  \item $\sigma(k) \geq \tau(k)$ para todo $k$
  \item $|a(\sigma(k)) - a(\tau(k))| \geq \varepsilon$ para todo $k$
\end{itemize}

Por la monotonicidad de $a$ y el orden $\tau(k) \leq \sigma(k)$, tenemos
$a(\tau(k)) \leq a(\sigma(k))$, luego:
\[
  \varepsilon \leq |a(\sigma(k)) - a(\tau(k))| = a(\sigma(k)) - a(\tau(k))
\]

Ahora aplicamos el teorema \texttt{IterateGap} para concluir:
\[
  \forall k \in \mathbb{N},\quad k \cdot \varepsilon \leq a(\sigma^{[k]}(0)) - a(0)
\]

Tomemos $k = \lceil(M - a(0))/\varepsilon\rceil_+ + 1$. Entonces:
\[
  k > \frac{M - a(0)}{\varepsilon}
\]
lo que implica:
\[
  k \cdot \varepsilon > M - a(0)
\]

De nuestro resultado de acumulación de brechas:
\[
  k \cdot \varepsilon \leq a(\sigma^{[k]}(0)) - a(0)
\]

Por tanto:
\[
  M - a(0) < k \cdot \varepsilon \leq a(\sigma^{[k]}(0)) - a(0)
\]

Esto nos da $M < a(\sigma^{[k]}(0))$, contradiciendo nuestra hipótesis de
que $a(n) \leq M$ para todo $n$.

Por lo tanto, $a$ debe ser de Cauchy.
\end{proof}

\section{Brechas iteradas: el lema auxiliar}

Ahora demostramos el resultado pendiente de la sección anterior. Este
lema técnico captura el mecanismo preciso por el cual las brechas
persistentes en una sucesión monótona se acumulan bajo iteración.

Recordemos nuestro contexto: tenemos una sucesión monótona
$a : \mathbb{N} \to X$ y dos subsucesiones $\sigma, \tau : \mathbb{N} \to
\mathbb{N}$ con las siguientes propiedades:
\begin{itemize}[nosep]
  \item $\sigma$ crece más rápido que $\tau$: $\sigma(n) \geq \tau(n)$
    para todo $n$
  \item $\tau$ crece más rápido que la función identidad: $\tau(n) \geq n$
    para todo $n$
  \item Existe algún $\varepsilon > 0$ tal que:
    $\varepsilon \leq |a(\sigma(n)) - a(\tau(n))|$
\end{itemize}

Nuestro objetivo es mostrar que si iteramos $\sigma$ exactamente $k$
veces --- escrito en Lean como $\sigma^{[k]}$ --- entonces acumularemos al
menos $k \cdot \varepsilon$ de crecimiento desde el valor inicial.

\subsection{El enunciado matemático}

\begin{theorem}[\texttt{IterateGap}]
Dada una sucesión monótona $a : \mathbb{N} \to X$ y subsucesiones
$\tau, \sigma : \mathbb{N} \to \mathbb{N}$ que satisfacen las condiciones
anteriores, se tiene:
\[
  \forall k \in \mathbb{N},\quad k \cdot \varepsilon \leq a(\sigma^{[k]}(0)) - a(0)
\]
\end{theorem}

\subsection{Estrategia de la demostración}

La demostración procede por inducción sobre $k$. La idea clave es que
cada iteración de $\sigma$ introduce al menos una brecha adicional de
tamaño $\varepsilon$, y la monotonicidad garantiza que estas brechas se
acumulan de manera aditiva.

Para el caso base $k = 0$, el enunciado se convierte en
$0 \leq a(0) - a(0) = 0$, que es trivial.

Para el paso inductivo, usamos la condición de brecha en el punto
$\sigma^{[k]}(0)$ para establecer que el movimiento de $\sigma^{[k]}(0)$
a $\sigma^{[k+1]}(0)$ crea una brecha adicional de al menos $\varepsilon$,
que se combina con la hipótesis inductiva.

\subsection{Demostración en Lean}

\begin{minted}{lean4}
Statement {X : Type*} [NormedField X] [LinearOrder X]
    [IsStrictOrderedRing X] (a : ℕ → X) (ha : Monotone a)
    (ε : X) (εpos : ε > 0) (τ : ℕ → ℕ) (hτ : ∀ n, τ n ≥ n)
    (σ : ℕ → ℕ) (hσ : ∀ n, σ n ≥ τ n)
    (hgap : ∀ n, ε ≤ |a (σ n) - a (τ n)|)
    : ∀ (k : ℕ), k * ε ≤ a (σ^[k] 0) - a 0 := by
  intro k
  induction' k with k hk
  · norm_num
  · specialize hgap (σ^[k] 0)
    rewrite [show |a (σ (σ^[k] 0)) - a (τ (σ^[k] 0))|
               = a (σ (σ^[k] 0)) - a (τ (σ^[k] 0))
             by apply abs_of_nonneg (by bound)] at hgap
    rewrite [show σ (σ^[k] 0) = σ^[k+1] 0 by apply succ_iterate] at hgap
    have f1 : 0 ≤ a (τ (σ^[k] 0)) - a (σ^[k] 0) := by bound
    push_cast
    linarith [f1, hk, hgap]
\end{minted}

\subsection{Demostración en lenguaje natural}

\begin{proof}
Procedemos por inducción sobre $k$.

\textbf{Caso base:} Cuando $k = 0$, necesitamos mostrar
$0 \cdot \varepsilon \leq a(\sigma^{[0]}(0)) - a(0)$. Como
$\sigma^{[0]}(0) = 0$, esto se convierte en $0 \leq a(0) - a(0) = 0$, que
es inmediato.

\textbf{Paso inductivo:} Supongamos que el resultado se cumple para algún
$k \geq 0$, es decir,
$k \cdot \varepsilon \leq a(\sigma^{[k]}(0)) - a(0)$.
Debemos mostrar $(k+1) \cdot \varepsilon \leq a(\sigma^{[k+1]}(0)) - a(0)$.

Por la hipótesis de brecha aplicada en $n = \sigma^{[k]}(0)$:
\[
  \varepsilon \leq |a(\sigma(\sigma^{[k]}(0))) - a(\tau(\sigma^{[k]}(0)))|
\]

Como $\sigma^{[k]}(0) \geq 0$ y $h_\tau$ implica
$\tau(\sigma^{[k]}(0)) \geq \sigma^{[k]}(0)$, y $h_\sigma$ implica
$\sigma(\sigma^{[k]}(0)) \geq \tau(\sigma^{[k]}(0))$, la monotonicidad
nos da:
\[
  a(\sigma^{[k]}(0)) \leq a(\tau(\sigma^{[k]}(0))) \leq a(\sigma(\sigma^{[k]}(0)))
\]

Por lo tanto, el valor absoluto es igual a la diferencia:
\[
  \varepsilon \leq a(\sigma(\sigma^{[k]}(0))) - a(\tau(\sigma^{[k]}(0)))
\]

Usando \texttt{succ\_iterate}, tenemos
$\sigma(\sigma^{[k]}(0)) = \sigma^{[k+1]}(0)$.
Además, por monotonicidad:
\[
  a(\tau(\sigma^{[k]}(0))) - a(\sigma^{[k]}(0)) \geq 0
\]

Combinando estas desigualdades:
\begin{align*}
  a(\sigma^{[k+1]}(0)) - a(0)
    &= a(\sigma^{[k+1]}(0)) - a(\sigma^{[k]}(0)) + a(\sigma^{[k]}(0)) - a(0) \\
    &\geq a(\sigma^{[k+1]}(0)) - a(\tau(\sigma^{[k]}(0))) + a(\sigma^{[k]}(0)) - a(0) \\
    &\geq \varepsilon + k \cdot \varepsilon \\
    &= (k+1) \cdot \varepsilon
\end{align*}

Esto completa la inducción.
\end{proof}

\subsection{Por qué este lema técnico importa}

Este teorema captura el mecanismo preciso por el cual las sucesiones
monótonas con brechas persistentes eventualmente violan cualquier cota
superior propuesta. La iteración de subsucesiones, combinada con la
acumulación aditiva de brechas, proporciona la herramienta cuantitativa
necesaria para hacer riguroso el argumento intuitivo.

El resultado demuestra un principio fundamental: en sistemas monótonos,
los patrones de crecimiento local (representados por la condición de
brecha) escalan linealmente bajo iteración. Esta propiedad de escala es
lo que nos permite construir una contradicción con el acotamiento en el
teorema principal.

%%% Local Variables:
%%% mode: LaTeX
%%% TeX-master: "Analisis_en_Lean"
%%% End:
